\subsection{Embodied Agents}
\label{sec:app-embodied}


Embodied agents extend reasoning beyond text, anchoring language in robotic perception, manipulation and navigation.  By embedding LLMs within robotic and simulated bodies, these embodied agents tackle real-world generalization, continual adaptation and multi-modal grounding.  

In this subsection, we begin with the foundational layer (Section~\ref{sec:app-embodied-planning}), which covers long-horizon embodied planning, tool-assisted perception, manipulation and execution. Building upon these capabilities, the self-evolving layer (Section~\ref{sec:app-embodied-evolve}) introduces agentic memory, feedback and self-reflection capabilities  enabling robots to refine control policies, adapt to novel environments and improve performance through continual interaction. Finally, the collective reasoning layer explores multi-robot collaboration (Section~\ref{sec:app-embodied-multi}), where agents coordinate perception, share learned representations and jointly reason about tasks to achieve complex embodied goals.

\subsubsection{Foundational agentic reasoning}


\label{sec:app-embodied-planning}

\paragraph{Planning.}


Early work such as SayCan~\cite{ahn2022icanisay} established the template by mapping linguistic descriptions to skill affordance estimates and  SayPlan~\cite{DBLP:conf/corl/RanaHGA0S23} refined this grounding by leveraging 3D scene graphs to align goal references with object-centric representations and spatial models. Beyond symbolic representations, EmbodiedGPT~\cite{mu2023embodiedgptvisionlanguagepretrainingembodied} use curated video CoT annotations of sub-goals to train models taht map multi-model input to structured sequences for embodied planning, while context-aware planning system~\cite{kim2024contextawareplanningenvironmentawarememory} adds semantic spatial map and object location information to the planning pipeline, enabling dynamical planning during execution.  In addition, DEPS~\cite{DBLP:journals/corr/abs-2302-01560} introduces an interactive planning loop (i.e. describe, explain, plan and select) for open-world multi-task agents. 


Embodied agents also rely on multi-modal reasoning traces that explicitly align perception with action. For example, Embodied CoT~\cite{zawalski2024robotic} trains vision-language-action models to generate reasoning steps incorporating visual features before executing an action. Fast\,ECoT~\cite{duan2025fast} accelerates this by caching and re-using reasoning segments across time-steps, reducing inference latency while preserving task success. More recently, Cosmos-Reason1~\cite{nvidia2025cosmosreason1physicalcommonsense} establishes an ontology of space, time and dynamics that lets CoT sequences encode structured physical priors. CoT-VLA~\cite{zhao2025cot} builds a visual chain-of-thought by predicting future image frames as intermediate sub-goals prior to action generation. Finally, Emma-X~\cite{sun2024emmaxembodiedmultimodalaction} integrates grounded chain-of-thought with look-ahead spatial reasoning, improving long-horizon embodied task performance.


Another line of works strengthen embodied planning through reinforcement learning, considering planning not only as static decomposition but as a self-evolving process that adapts to environment feedback. Robot-R1~\cite{kim2025robotr1reinforcementlearningenhanced} trains large VLMs to predict keypoint transitions under visual context, turning RL into a mechanism for learning physically grounded forward models. ManipLVM-R1~\cite{song2025maniplvm} exploits verifiable physical reward signals (e.g., trajectory match and affordance correctness) to reduce reliance on dense expert annotation. Embodied-R~\cite{zhao2025embodied} presents a collaborative framework where VLMs handle perception and smaller LMs handle reasoning, and the whole is trained via RL for embodied spatial reasoning. VIKI-R~\cite{kang2025viki} further extends this direction into heterogeneous multi-agent cooperation with a two-stage design, employing a two‐stage pipeline of chain-of-thought fine-tuning followed by hierarchical RL across agents coordinating activation and planning. 



\paragraph{Tool-use.}

\label{sec:app-embodied-tool}


Embodied agents can also be strengthened to interact with external tools to enhance perception and compensate for incomplete observations. GSCE~\cite{wang2025gscepromptframeworkenhanced}, for example, provides a prompt-framework that binds skill APIs and constraints for safe LLM-driven drone control. MineDojo~\cite{fan2022minedojo} links agents to internet-scale corpora and thus enabling richer affordance grounding. Physical AI Agents~\cite{bousetouane2025physicalaiagentsintegrating} further introduces a modular architecture and a retrieval augmented generation design pattern for embedding real-world physical interaction into LLM-driven agents. Beyond offline tool use, some systems treat the environment itself as an API. For example, Matcha agent~\cite{zhao2023chatenvironmentinteractivemultimodal} uses an LLM to issue queries about objects and scenes and thereby acquire perceptual information needed for task completion.



On the other hand, execution module is one of the most important tool type. It translates high-level language instructions into continuous motor commands, enabling embodied agents to act reliably in physical environments.  Early systems such as SayCan~\cite{ahn2022icanisay} uses language to invoke robot pick-and-place skills; while LEO~\cite{huang2024embodiedgeneralistagent3d} broaden execution to more general manipulation settings and Hi Robot~\cite{shi2025hi} uses a VLM reasoner to process complex prompts and a low-level action policy executes the chosen step. More recent efforts broaden the execution space: Gemini Robotics~\cite{geminiroboticsteam2025geminiroboticsbringingai} introduces a large-scale vision-language-action model for real-world robot control and Octopus~\cite{yang2024octopusembodiedvisionlanguageprogrammer} generates executable code in simulated environments that bridges planning and manipulation. 


Beyond single-agent control, hybrid pipelines couple reactive reflexes with language-guided policies to support complex domains. For example, CaPo~\cite{liu2024capo} incorporates an execution phase where agents carry out decomposed sub-tasks and adapt their meta-plan based on progress; COHERENT~\cite{liu2024coherent} embeds a robot executor module within its PEFA (i.e. proposal, execution, feedback and adjustment) loop, which ensures each assigned sub-task is acted and refined appropriately; and MP5~\cite{qin2024mp5multimodalopenendedembodied} integrates multi-modal perception to generate executable plans in open-ended Minecraft. At the perception–action interface, LLM-Planner~\cite{song2023llmplannerfewshotgroundedplanning} generates sub-goals and maps them into action sequences via a low-level controller and EmbodiedGPT~\cite{mu2023embodiedgptvisionlanguagepretrainingembodied} illustrate how LLM-generated plans can be translated into control policy for embodied control in physical environments. 



\paragraph{Search and retrieval.}

Embodied agents can also use search and retrieval ability to ground language in spatial structure and past experience. Early navigation systems such as L3MVN~\cite{yu2023l3mvn} use LLMs to query a semantic map and select promising frontiers as long-term goals during visual target navigation, while SayNav~\cite{rajvanshi2024saynav} and SayPlan~\cite{DBLP:conf/corl/RanaHGA0S23} build 3D scene graphs and then search task-relevant subgraphs so language instructions can be translated into grounded waypoints and sub-tasks in large environments. Long-horizon navigation works like ReMEmbR~\cite{anwar2025remembr} maintain a structured spatio-temporal memory that can be queried to answer “where” and “when” questions about past robot experience. Additionally, RAG-style systems make retrieval a first-class part of the planning loop: Embodied-RAG~\cite{quanting2024embodiedrag} and EmbodiedRAG~\cite{meghan2024embodiedrag} treat an agent’s experience and 3D scene graphs as non-parametric memories from which task-relevant episodes or subgraphs are retrieved for navigation and task planning; Retrieval-Augmented Embodied Agents~\cite{zhu2024retrieval} retrieve policies from a shared memory bank and condition action generation on them; and MLLM-as-Retriever~\cite{junpeng2024mllm} trains a multi-modal LLM retriever to rank past trajectories so each decision step can condition on the most useful prior experience rather than only the current observation.


\subsubsection{Self-evolving agentic reasoning}
\label{sec:app-embodied-evolve}



Embodied agents reliably achieve long-horizon autonomy when they can self-evolve over time: monitor their own internal states, store and update task-relevant knowledge and adjust behaviors when plans deviate. In the following paragraphs, we examine how memory modules, feedback signals and agentic reflection enable embodied agents to turn planning from a one-shot process into a continually improving cycle of behavior.


\paragraph{Memory.}
Effective memory mechanisms enable agents to reuse past experiences and maintain coherent task execution over extended interactions. Many systems cache recent observations in episodic buffers while summarizing long-term semantics in structured graphs, as in household planning~\cite{glocker2025llm} and long-horizon agents with hybrid multi-modal memory~\cite{li2024optimus1}. Skills and routines can be shared across tasks via indexed memory stores. For example, HELPER-X~\cite{sarch2023openendedinstructableembodiedagents} indexes discovered skills and action scripts, which aid future dialogue and can be shared across domains.
Spatial navigation methods such as BrainNav~\cite{ling2025endowingembodiedagentsspatial} maintain biologically inspired dual-map memories linked by a hippocampal hub to reduce hallucinations and drift. Broader contexts also benefit: CAPEAM~\cite{kim2024contextawareplanningenvironmentawarememory} incorporates environment-aware memory modules that track object states and spatial changes. Finally, lifelong episodic systems such as Ella~\cite{zhang2025ellaembodiedsocialagents} maintains long-term multi-modal memory system to support social-robot interaction.


\paragraph{Agentic Feedback and Reflection.}
Dialogue-based critique, calibrated uncertainty and environment-aware reward shaping refine policies beyond binary success signals. For example, Matcha agent~\cite{zhao2023chatenvironmentinteractivemultimodal} treats objects and scenes as interactive information sources before acting and FAMER~\cite{wang2025communicationefficientdesirealignmentembodied} uses lightweight preference feedback to adapt embodied agents to user intentions in real time. Uncertainty-aware planners such as KnowNo~\cite{ren2023robotsaskhelpuncertainty}, which proactively solicit guidance when confidence falls below guarantees, and Octopus~\cite{yang2024octopusembodiedvisionlanguageprogrammer}, which exploits environmental feedback to improve generated executable programs over time. At the multi-agent level, MindForge~\cite{licua2024mindforge} introduces theory-of-mind style perspective feedback so heterogeneous robots adapt to each other’s reasoning strategies; while ReAd~\cite{zhang2024efficientllmgroundingembodied} introduces a advantage-based feedback loop that enables an LLM planner to self-refine its collaboration strategies across embodied multi-agent tasks.



Robust reflection mechanisms help agents anticipate failures by monitoring their own reasoning and actions and then adjusting plans. Optimus-1~\cite{li2024optimus1} couples a \emph{Knowledge-guided Planner} with an \emph{Experience-Driven Reflector} to revise decisions using stored experience, while another recent study~\cite{Moncada_Ramirez_2025} defines structured agentic workflows (including self-Reflection, multi-Agent reflection and LLM Ensemble) that enable robots to reflect on and refine LLM-generated object-centered plans, thus reducing reasoning errors. Systems such as EMAC+~\cite{ao2025emac} interleave perception, planning and verification steps to perform online plan refinement and earlier works such as Voyager~\cite{wang2023voyager} also embeds an iterative prompting loop that uses environment feedback and execution errors to refine its skill library over time.


\subsubsection{Collective multi-agent reasoning}
\label{sec:app-embodied-multi}


Multi-agent collaboration enables embodied systems to divide labor and coordinate complex tasks more efficiently, with language often serving as the primary medium for negotiation and role allocation. For instance, SMART-LLM~\cite{kannan2024smartllmsmartmultiagentrobot} decomposes high-level instructions and allocates sub-tasks across multiple robots, while CaPo~\cite{liu2024capo} optimizes cooperative plans to avoid redundant exploration. For heterogeneity and coordination mechanisms, COHERENT~\cite{liu2024coherent} deploys a propose-execution-feedback-adjust loop across diverse robot types to enable seamless joint operation. In addition, Theory of Mind (ToM), which refers to an embodied agent’s ability to infer and reason about others' beliefs and mental states, is also highly related to embodied multi-agent systems~\cite{li2023theory, wu2025large, cross2024hypothetical}. For example, MindForge~\cite{licua2024mindforge} equips agents with explicit theory-of-mind representations and natural inter-agent communication to coordinate collaboratively.


For multi-modal frameworks, EMAC+~\cite{ao2025emac} integrate vision and language modules and continuously refine plans via visual feedback, COMBO~\cite{zhang2025combocompositionalworldmodels} integrates vision and language modules and continuously refine plans via visual feedback, and VIKI-R~\cite{kang2025viki} demonstrates reinforcement learning as a scalable coordination mechanism among embodied agents. At larger scales, studies such as RoCo~\cite{mandi2024roco} show how role negotiation and flexible protocols support adaptable teamwork in dynamic environments.
